\textbf{Phần 1. Trắc nghiệm lựa chọn}
\setcounterpageref{ex}{0}

\begin{ex} %[5D2N14]
Số thập phân gồm ``Ba chục, bảy đơn vị, ba phần mười, sáu phần trăm, bốn phần nghìn'' được viết là
\choice
{$307{,}364$}
{$37{,}346$}
{$373{,}64$}
{\True $37{,}364$}
\loigiai{
Phần nguyên gồm $3$ chục và $7$ đơn vị là $37$. Phần thập phân gồm $3$ phần mười, $6$ phần trăm, $4$ phần nghìn là $364$. Số đó là $37{,}364$.\\
\textbf{Đáp án của câu hỏi là D}
}
\end{ex}

\begin{ex} %[5D2N14]
Số thập phân gồm ``Không đơn vị, hai mươi lăm phần nghìn'' được viết là
\choice
{$0{,}25$}
{$0{,}250$}
{\True $0{,}025$}
{$25{,}000$}
\loigiai{
Số có phần nguyên là $0$, hàng phần mười là $0$, hàng phần trăm là $2$, hàng phần nghìn là $5$ được viết là $0{,}025$.\\
\textbf{Đáp án của câu hỏi là C}
}
\end{ex}

\begin{ex} %[5D2H14]
Số thập phân gồm ``Một nghìn đơn vị, bảy mươi mốt phần nghìn'' được viết là
\choice
{$1\,000{,}71$}
{\True $1\,000{,}071$}
{$1\,071{,}0$}
{$1\,000{,}710$}
\loigiai{
Phần nguyên là $1\,000$, phần thập phân có hàng phần mười là $0$, hàng phần trăm là $7$, hàng phần nghìn là $1$ nên viết là $1\,000{,}071$.\\
\textbf{Đáp án của câu hỏi là B}
}
\end{ex}

\begin{ex} %[5D2H14]
Viết số đo $4\text{ cm } 6\text{ mm}$ dưới dạng số thập phân có đơn vị xăng-ti-mét ta được
\choice
{\True $4{,}6\text{ cm}$}
{$4{,}06\text{ cm}$}
{$46{,}0\text{ cm}$}
{$0{,}46\text{ cm}$}
\loigiai{
Ta có: $4\text{ cm } 6\text{ mm} = 4\dfrac{6}{10}\text{ cm} = 4{,}6\text{ cm}$.\\
\textbf{Đáp án của câu hỏi là A}
}
\end{ex}

\begin{ex} %[5D2H14]
Đổi $4{,}75\text{ kg}$ ra gam ta được kết quả là
\choice
{$475\text{ g}$}
{$4\,075\text{ g}$}
{$47\,500\text{ g}$}
{\True $4\,750\text{ g}$}
\loigiai{
Ta có: $4{,}75\text{ kg} = 4\text{ kg} + 0{,}75\text{ kg} = 4\,000\text{ g} + 750\text{ g} = 4\,750\text{ g}$.\\
\textbf{Đáp án của câu hỏi là D}
}
\end{ex}

\begin{ex} %[5D2H14]
Làm tròn số thập phân $13{,}56$ đến số tự nhiên gần nhất ta được
\choice
{\True $14$}
{$13$}
{$13{,}5$}
{$13{,}6$}
\loigiai{
Chữ số hàng phần mười là $5$ nên làm tròn lên thành $14$.\\
\textbf{Đáp án của câu hỏi là A}
}
\end{ex}

\begin{ex} %[5D2H14]
Làm tròn số thập phân $1{,}36$ đến hàng phần mười ta được
\choice
{$1{,}3$}
{$1{,}0$}
{\True $1{,}4$}
{$2{,}0$}
\loigiai{
Chữ số hàng phần trăm là $6 > 5$ nên làm tròn lên thành $1{,}4$.\\
\textbf{Đáp án của câu hỏi là C}
}
\end{ex}

\begin{ex} %[5D2H14]
Làm tròn số thập phân $1{,}026$ đến hàng phần trăm ta được kết quả là
\choice
{$1{,}02$}
{\True $1{,}03$}
{$1{,}2$}
{$1{,}00$}
\loigiai{
Chữ số hàng phần nghìn là $6 > 5$ nên làm tròn lên thành $1{,}03$.\\
\textbf{Đáp án của câu hỏi là B}
}
\end{ex}

\begin{ex} %[5D2H14]
Sắp xếp các số $3{,}527;\; 3{,}571;\; 2{,}752;\; 3{,}275$ theo thứ tự từ bé đến lớn là
\choice
{$3{,}571;\; 3{,}527;\; 3{,}275;\; 2{,}752$}
{$2{,}752;\; 3{,}527;\; 3{,}275;\; 3{,}571$}
{\True $2{,}752;\; 3{,}275;\; 3{,}527;\; 3{,}571$}
{$3{,}275;\; 2{,}752;\; 3{,}527;\; 3{,}571$}
\loigiai{
Ta có: $2{,}752 < 3{,}275 < 3{,}527 < 3{,}571$.\\
\textbf{Đáp án của câu hỏi là C}
}
\end{ex}

\begin{ex} %[5D2H14]
Ba ô tô chở hàng: xe A chở $2{,}546\text{ tấn}$, xe B chở $2{,}645\text{ tấn}$, xe C chở $2{,}564\text{ tấn}$. Xe chở lượng hàng nhẹ nhất là
\choice
{Xe B}
{\True Xe A}
{Xe C}
{Cả ba xe chở bằng nhau}
\loigiai{
So sánh: $2{,}546 < 2{,}564 < 2{,}645$. Do đó xe A chở lượng hàng nhẹ nhất.\\
\textbf{Đáp án của câu hỏi là B}
}
\end{ex}

\begin{ex} %[5D2V14]
Biết $1\text{ l}$ dầu ăn cân nặng $0{,}9\text{ kg}$, $1\text{ l}$ mật ong cân nặng $1{,}36\text{ kg}$. Hỏi $2\text{ l}$ dầu ăn nhẹ hơn $2\text{ l}$ mật ong bao nhiêu ki-lô-gam?
\choice
{\True $0{,}92\text{ kg}$}
{$0{,}46\text{ kg}$}
{$1{,}02\text{ kg}$}
{$0{,}82\text{ kg}$}
\loigiai{
$1\text{ l}$ dầu ăn nhẹ hơn $1\text{ l}$ mật ong là: $1{,}36 - 0{,}9 = 0{,}46\text{ (kg)}$.\\
$2\text{ l}$ dầu ăn nhẹ hơn $2\text{ l}$ mật ong là: $0{,}46 \times 2 = 0{,}92\text{ (kg)}$.\\
\textbf{Đáp án của câu hỏi là A}
}
\end{ex}

\begin{ex} %[5D2V14]
Từ các thẻ số $5$; $1$; $7$; $0$ và dấu phẩy ``$,$'', số thập phân lớn nhất bé hơn $1$ (mỗi thẻ dùng đúng một lần) là
\choice
{$0{,}715$}
{$0{,}571$}
{$1{,}057$}
{\True $0{,}751$}
\loigiai{
Để số bé hơn $1$, chữ số phần nguyên phải là $0$. Để số lớn nhất, các chữ số phần thập phân được sắp xếp giảm dần: $7, 5, 1$. Số đó là $0{,}751$.\\
\textbf{Đáp án của câu hỏi là D}
}
\end{ex}

\textbf{Phần 2. Trắc nghiệm đúng sai}
\setcounterpageref{ex}{0}

\begin{ex} %[5D2H14]
Cho số thập phân $A = 508{,}42$ và số thập phân $B = 37{,}364$. Xét tính đúng sai của các khẳng định sau:
\choiceTF
{\True Số $A$ được đọc là: Năm trăm linh tám phẩy bốn mươi hai.}
{Chữ số $4$ trong số $A$ và chữ số $4$ trong số $B$ có giá trị bằng nhau.}
{\True Làm tròn số $A$ đến số tự nhiên gần nhất ta được $508$.}
{\True Làm tròn số $B$ đến hàng phần mười ta được $37{,}4$.}
\loigiai{
\begin{itemize}
    \item Số $A = 508{,}42$ đọc là: Năm trăm linh tám phẩy bốn mươi hai.
    \item Chữ số $4$ trong số $A$ ở hàng phần mười (giá trị $\dfrac{4}{10}$), chữ số $4$ trong số $B$ ở hàng phần nghìn (giá trị $\dfrac{4}{1\,000}$), hai giá trị này khác nhau.
    \item Số $A$ có hàng phần mười là $4 < 5 \Rightarrow$ làm tròn thành $508$.
    \item Số $B$ có hàng phần trăm là $6 > 5 \Rightarrow$ làm tròn đến hàng phần mười là $37{,}4$.
\end{itemize}
\begin{itemize}
    \item a) \textbf{Đúng}.
    \item b) \textbf{Sai}: Giá trị của chúng khác nhau.
    \item c) \textbf{Đúng}.
    \item d) \textbf{Đúng}.
\end{itemize}
}
\end{ex}

\begin{ex} %[5D2V14]
Khối lượng của $1\text{ l}$ một số chất lỏng được cho như sau: Nước biển ($1{,}026\text{ kg}$), Rượu ($0{,}79\text{ kg}$), Mật ong ($1{,}36\text{ kg}$), Thuỷ ngân ($13{,}56\text{ kg}$). Xét tính đúng sai của các khẳng định sau:
\choiceTF
{\True Làm tròn khối lượng $1\text{ l}$ nước biển đến hàng phần mười ta được $1{,}0\text{ kg}$ (hay $1\text{ kg}$).}
{Làm tròn khối lượng $1\text{ l}$ rượu đến số tự nhiên gần nhất ta được $0\text{ kg}$.}
{\True Làm tròn khối lượng $1\text{ l}$ mật ong đến hàng phần mười ta được $1{,}4\text{ kg}$.}
{\True Trong các chất lỏng trên, thuỷ ngân có khối lượng nặng nhất.}
\loigiai{
\begin{itemize}
    \item Nước biển: $1{,}026$ có hàng phần trăm là $2 < 5 \Rightarrow$ làm tròn đến hàng phần mười là $1{,}0\text{ kg}$.
    \item Rượu: $0{,}79$ có hàng phần mười là $7 > 5 \Rightarrow$ làm tròn đến số tự nhiên gần nhất là $1\text{ kg} \neq 0\text{ kg}$.
    \item Mật ong: $1{,}36$ có hàng phần trăm là $6 > 5 \Rightarrow$ làm tròn đến hàng phần mười là $1{,}4\text{ kg}$.
    \item Thuỷ ngân có khối lượng $13{,}56\text{ kg}$, lớn nhất trong các chất lỏng trên.
\end{itemize}
\begin{itemize}
    \item a) \textbf{Đúng}.
    \item b) \textbf{Sai}: Làm tròn thành $1\text{ kg}$.
    \item c) \textbf{Đúng}.
    \item d) \textbf{Đúng}.
\end{itemize}
}
\end{ex}

\textbf{Phần 3. Trả lời ngắn}
\setcounterpageref{ex}{0}

\begin{ex} %[5D2H14]
Viết số đo $5\text{ kg } 256\text{ g}$ dưới dạng số thập phân có đơn vị ki-lô-gam.
\shortans{5,256}
\loigiai{
\textbf{Đáp số:} $5{,}256$\\
$5\text{ kg } 256\text{ g} = 5\dfrac{256}{1\,000}\text{ kg} = 5{,}256\text{ kg}$.
}
\end{ex}

\begin{ex} %[5D2H14]
Đổi $2{,}35\text{ m}$ ra xăng-ti-mét. Kết quả là bao nhiêu xăng-ti-mét?
\shortans{235}
\loigiai{
\textbf{Đáp số:} $235$\\
$2{,}35\text{ m} = 235\text{ cm}$.
}
\end{ex}

\begin{ex} %[5D2H14]
Làm tròn số thập phân $0{,}07085$ đến hàng phần trăm. (Nhập kết quả dưới dạng số thập phân).
\shortans{0,07}
\loigiai{
\textbf{Đáp số:} $0{,}07$\\
Chữ số ở hàng phần nghìn là $0 < 5$ nên làm tròn thành $0{,}07$.
}
\end{ex}

\begin{ex} %[5D2H14]
Tìm số tự nhiên $x$ lớn nhất thỏa mãn: $x < 37{,}364$.
\shortans{37}
\loigiai{
\textbf{Đáp số:} $37$\\
Số tự nhiên lớn nhất nhỏ hơn $37{,}364$ là $37$.
}
\end{ex}

\begin{ex} %[5D2V14]
Từ các thẻ số $5$; $1$; $7$; $0$ và dấu phẩy ``$,$'', lập số thập phân bé nhất lớn hơn $0$ và bé hơn $1$ (mỗi thẻ dùng đúng một lần). Hỏi số đó là số nào? (Nhập dưới dạng số thập phân có dấu phẩy).
\shortans{0,157}
\loigiai{
\textbf{Đáp số:} $0{,}157$\\
Phần nguyên là $0$. Để số bé nhất, các chữ số phần thập phân sắp xếp tăng dần: $1, 5, 7$. Số đó là $0{,}157$.
}
\end{ex}

\begin{ex} %[5D2V14]
Một xe ô tô chở $2{,}546\text{ tấn}$ hàng, một xe khác chở $2{,}645\text{ tấn}$ hàng. Hỏi cả hai xe chở tất cả bao nhiêu ki-lô-gam hàng?
\shortans{5191}
\loigiai{
\textbf{Đáp số:} $5191$\\
Đổi $2{,}546\text{ tấn} = 2\,546\text{ kg}$; $2{,}645\text{ tấn} = 2\,645\text{ kg}$.\\
Tổng khối lượng hàng là: $2\,546 + 2\,645 = 5\,191\text{ (kg)}$.
}
\end{ex}
